\section{请求体的落盘与调试}

\subsection{如何查看完整请求体}

GPT 的 \texttt{response.create} API 默认不落盘请求体。但 Codex 提供了调试选项：

\begin{lstlisting}[caption={启用请求体落盘的环境变量}, language=bash]
export CODEX_LOG_REQUESTS=1
# 请求体将保存到 ~/.codex/logs-1/skill-lite
\end{lstlisting}

落盘的 JSON 文件包含完整的请求体：instructions、developer 消息、user 消息、12 个 Tool 定义（含 Apply Patch 的 Lark Grammar）等。

\subsection{本章小结}

通过环境变量启用落盘，可以完整审视 Codex 发送给 API 的请求体，是调试和理解 Agent 行为的重要手段。


